%% ARITH 2027 Submission Scaffold
%% =================================
%% IEEEtran Conference class, two-column.
%% Status: v0.1 SCAFFOLD, 2026-06-01.
%% This file is the reviewer-facing build target. Anonymised per
%% arith2027_anonymisation_checklist.md. At camera-ready, replace
%% the \author{} block, restore \thanks{}, and re-enable the
%% acknowledgements section.

\documentclass[conference]{IEEEtran}
\IEEEoverridecommandlockouts

\usepackage{cite}
\usepackage{amsmath,amssymb,amsfonts}
\usepackage{algorithmic}
\usepackage{graphicx}
\usepackage{textcomp}
\usepackage{xcolor}
\usepackage{booktabs}
\usepackage{hyperref}

%% ASCII-only headers per goldenfloat-ladder skill rules.
%% phi is rendered as \varphi in math mode.

\def\BibTeX{{\rm B\kern-.05em{\sc i\kern-.025em b}\kern-.08em
    T\kern-.1667em\lower.7ex\hbox{E}\kern-.125emX}}

\begin{document}

\title{A Phi-Derived Static-Split Floating-Point Family\\
       from GF4 to GF256 with a Lucas-Exact Integer Identity}

%% --- ANONYMISED AUTHOR BLOCK FOR DOUBLE-BLIND REVIEW ---
\author{
\IEEEauthorblockN{Anonymous Submission}
\IEEEauthorblockA{ARITH 2027 --- Double-Blind Review}
}
%% --- END ANONYMISED BLOCK ---
%%
%% CAMERA-READY (DO NOT UNCOMMENT FOR REVIEWER PDF):
%% \author{
%% \IEEEauthorblockN{Dmitrii Vasiliev}
%% \IEEEauthorblockA{\textit{Independent Researcher} \\
%% Ko Samui, Thailand \\
%% admin@t27.ai \\
%% ORCID: 0009-0008-4294-6159}
%% }

\maketitle

\begin{abstract}
We describe GoldenFloat (GF), a family of floating-point formats whose
exponent:mantissa split is generated, for each total width $N$, by a
single closed rule $e = \mathrm{round}((N-1)/\varphi^2)$ with mantissa
$m = N-1-e$, where $\varphi = (1+\sqrt{5})/2$. The rule reproduces the
realised exponent widths of nine GF formats (GF4 through GF256) exactly
[Verified] and extends consistently to GF128, GF512, and GF1024. We
verify a classical integer identity, $\varphi^{2n} + \varphi^{-2n} =
L_{2n}$ (Lucas numbers), symbolically for $n = 1\ldots 256$ and
numerically at 500 decimal places, with classical attribution to Lucas
(1878) and the Binet formula. We report one hardware artefact: a GF16
codec on Artix-7 FPGA passing a 35-of-35 testbench at 323 MHz, and a
pre-registered matched-substrate head-to-head experiment against
posit-16, takum-16, and binary16 on the same board. We are explicit
about non-claims: $\varphi$ is not a uniqueness claim, no GF rung is
claimed to outperform posit, takum, MX, or LNS at matched substrate, and
no silicon has been validated. A falsification ledger records the open
questions and the experiments that would settle them.
\end{abstract}

\begin{IEEEkeywords}
computer arithmetic, floating-point formats, golden ratio, Lucas
numbers, FPGA, double-blind submission
\end{IEEEkeywords}

%% ===================== SECTION 1 =====================
\section{Introduction}
\label{sec:intro}

%% TODO: Fill from arith2027_paper_outline.md Section 1 (~525 words, 0.75 pg).
%% Three subsections:
%%   1. Motivation paragraph (posit, takum, MX, LNS as prior art).
%%   2. The three-word contribution: rule, identity, measurement.
%%   3. What this paper does not claim (four condensed items).

\subsection{Motivation}
% PLACEHOLDER: posit (Gustafson-Yonemoto 2017), takum (Hunhold ARITH
% 2025), OCP MX (Rouhani et al. 2023), LNS. Prior art and allies.

\subsection{Contribution}
% PLACEHOLDER: rule, identity, measurement.

\subsection{What This Paper Does Not Claim}
% PLACEHOLDER: (1) phi not unique; (2) no per-rung superiority;
% (3) no silicon validation; (4) breadth-as-moat is [Open-conjecture].

%% ===================== SECTION 2 =====================
\section{The Ladder Rule and Look-Elsewhere Correction}
\label{sec:ladder}

%% TODO: Fill from preprint Sections 2.1-2.4 (~1050 words, 1.5 pg).
%% Preserve the "negative result first" structure of the look-elsewhere
%% correction --- it is the strongest reviewer-signal subsection.

\subsection{Definition}
Let $N$ be the total format width in bits, including one sign bit. The
GoldenFloat split allocates
\begin{equation}
\label{eq:rule}
e = \mathrm{round}\!\left(\frac{N-1}{\varphi^2}\right), \qquad m = N - 1 - e,
\end{equation}
where $\varphi = (1+\sqrt{5})/2 = 1.61803\ldots$ and $\varphi^2 =
\varphi + 1 = 2.61803\ldots$, so $1/\varphi^2 = 0.38197\ldots$ The rule
is \emph{static}: $e$ depends only on $N$, never on the value being
represented.

% PLACEHOLDER: Table 1 with 12 formats (N, e, m, raw, ratio).

\subsection{Retraction of the ``2-to-256'' Wording}
% PLACEHOLDER: GF2 and GF3 are [Retracted]; ladder defined GF4 to GF256.

\subsection{Look-Elsewhere Correction for the $\varphi^{-2}$ Coincidence}
% PLACEHOLDER: 83-of-80000 rational search, family-wise probability
% approx 7.1e-3, Bonferroni saturates at 1, Bayes factor 964 with the
% structural caveat, narrowing from 392 to 47 ratios with 12 formats.

\subsection{Rounding Mode}
% PLACEHOLDER: round-half-to-even; immaterial for all nine [Verified]
% widths because the raw value is never a half-integer.

%% ===================== SECTION 3 =====================
\section{The GRE Multiplier-Free Anchor}
\label{sec:gre}

%% TODO: Fill from preprint Section 3 (~525 words, 0.75 pg).
%% Daubechies et al. IEEE TIT 2010 + the narrow defensible link.

\subsection{The Engineering Property}
% PLACEHOLDER: phi^2 = phi + 1 removes the constant-multiplier component
% in the beta-encoder of Daubechies, Gunturk, Wang, Yilmaz.

\subsection{What the GRE Result Does and Does Not Establish}
% PLACEHOLDER: (a) no claim base is forced; (b) no point-robustness at
% beta = phi (Daubechies Theorem 8 caveat); (c) the narrow defensible
% link to the Lucas identity of Section 4.

%% ===================== SECTION 4 =====================
\section{The Lucas-Exact Integer Identity}
\label{sec:lucas}

%% TODO: Fill from preprint Section 4 (~525 words, 0.75 pg).
%% Classical Lucas L_{2n} attribution + symbolic + 500-digit numerical
%% verification + application to GF bias derivation.

\subsection{Statement and Classical Provenance}
For all integers $n \geq 1$,
\begin{equation}
\label{eq:lucas}
\varphi^{2n} + \varphi^{-2n} = L_{2n},
\end{equation}
where $L_k$ is the $k$-th Lucas number ($L_0 = 2$, $L_1 = 1$, $L_k =
L_{k-1} + L_{k-2}$). This is a classical corollary of the Binet
formula $L_k = \varphi^k + (-\varphi)^{-k}$. The identity is
attributed to Lucas (1878) and Binet; the present authors claim no
originality for the identity itself.

\subsection{Numerical Verification}
% PLACEHOLDER: symbolic verification for n = 1..256; 500-digit mpmath
% numerical check; reproducibility script in supplementary materials.

\subsection{Application to GoldenFloat Bias Derivation}
% PLACEHOLDER: integer-backed accumulation idea; the engineering
% artefact does not yet exist [Open].

%% ===================== SECTION 5 =====================
\section{Hardware Status and the Matched-Substrate Experiment}
\label{sec:hardware}

%% TODO: Fill from preprint Section 5 + arith2027_matched_substrate_experiment.md.
%% This is the section that materially extends the preprint and makes
%% the ARITH paper viable.

\subsection{Per-Width Status}
% PLACEHOLDER: compressed Table 2; seven widths generated end-to-end
% from one .t27 spec into C + Verilog + JSON conformance vectors
% [Verified].

\subsection{The GF16 Result}
We report a GF16 codec implementation on Artix-7 FPGA passing a 35-of-35
testbench at 323 MHz after a one-line rounding fix [Verified]. Synthesis
was performed in Xilinx Vivado with default Performance\_Explore
implementation strategy. % PLACEHOLDER: detailed numbers + reference to
% supplementary RTL.

\subsection{The Substrate-Confound Problem}
% PLACEHOLDER: why software BPB measurements alone are insufficient.

\subsection{Matched-Substrate Head-to-Head Experiment}
\label{subsec:matched-substrate}
We pre-register a matched-substrate head-to-head experiment comparing
GF16 against posit-16 (PERCIVAL reference), takum-16 (open-source
reference), and binary16 (IEEE 754) on the same Artix-7 board.
% PLACEHOLDER: full protocol from arith2027_matched_substrate_experiment.md.
% Reports the four metrics across three workloads. All four outcomes
% (advantage, coexistence, disadvantage, insufficient-evidence) are
% acceptable; the experiment is designed to be falsifiable.

\subsection{Implications}
% PLACEHOLDER: how the matched-substrate result (or its pre-registration
% alone) resolves the substrate-confound; status update on the
% breadth-as-moat hypothesis.

%% ===================== SECTION 6 =====================
\section{Prior Art and Positioning}
\label{sec:prior}

%% TODO: Fill from preprint Section 6 (~350 words, 0.5 pg).

%% PLACEHOLDER: posit (Gustafson-Yonemoto 2017), takum (Hunhold 2024,
%% ARITH 2025) as closest live alternative, OCP MX (Rouhani et al.
%% 2023), LNS. Explicit: no accuracy or superiority claim.

%% ===================== SECTION 7 =====================
\section{Discussion: Breadth as the Moat}
\label{sec:moat}

%% TODO: Fill from preprint discussion + goldenfloat-ladder skill rules
%% (~350 words, 0.5 pg).

%% PLACEHOLDER: honest moat framing --- breadth (one rule, GF4 to GF256)
%% and toolchain coherence (one .t27 spec to C + Verilog + JSON across
%% seven widths). Not per-rung superiority. F1/F2/F3 falsification
%% protocol. Status: [Open-conjecture].

%% ===================== SECTION 8 =====================
\section{Conclusions and Falsification Path}
\label{sec:conc}

%% TODO: Fill from preprint conclusion (~175 words, 0.25 pg).

%% PLACEHOLDER: restate rule, identity, measurement contributions;
%% restate non-claims; pointer to FL-002 ledger in supplementary
%% materials; one sentence on next steps.

%% ===================== ACKNOWLEDGEMENTS =====================
%% REVIEWER-FACING: removed entirely for double-blind.
%%
%% CAMERA-READY:
%% \section*{Acknowledgements}
%% The author thanks % TODO --- check independence rules before listing.

%% ===================== REFERENCES =====================
\begin{thebibliography}{99}

\bibitem{daubechies2010gre}
I. Daubechies, C. S. G\"unt\"urk, Y. Wang, and \"O. Yilmaz, ``The Golden
Ratio Encoder,'' \emph{IEEE Transactions on Information Theory},
vol.~56, no.~10, pp.~5097--5110, Oct.\ 2010.

\bibitem{gustafson2017posit}
J.~L. Gustafson and I.~Yonemoto, ``Beating Floating Point at its Own
Game: Posit Arithmetic,'' \emph{Supercomputing Frontiers and
Innovations}, vol.~4, no.~2, 2017.

\bibitem{hunhold2024takum}
L.~Hunhold, ``Takum Arithmetic: A Tapered Floating-Point Format,''
arXiv:2412.20273, 2024.

\bibitem{hunhold2025fft}
L.~Hunhold and J.~L. Gustafson, ``Numerical analysis of FFT using takum
arithmetic,'' arXiv:2504.21197, 2025.

\bibitem{hunhold2025arnoldi}
L.~Hunhold and J.~L. Gustafson, ``An Arnoldi-style method for takum
arithmetic,'' arXiv:2504.21130, 2025.

\bibitem{rouhani2023mx}
B.~Rouhani \emph{et al.}, ``Microscaling Data Formats for Deep
Learning,'' OCP MX White Paper, 2023.

\bibitem{mallasen2022percival}
R.~Murillo, A.~A. Del Barrio, G.~Botella, M.~S. Kim, H.~Kim, and N.~Bagherzadeh,
``PERCIVAL: Open-Source Posit RISC-V Core with Quire Capability,''
\emph{IEEE Transactions on Emerging Topics in Computing}, 2022.

\bibitem{lucas1878}
E.~Lucas, ``Th\'eorie des fonctions num\'eriques simplement
p\'eriodiques,'' \emph{American Journal of Mathematics}, vol.~1, 1878.

\bibitem{ieee754-2019}
IEEE, ``IEEE Standard for Floating-Point Arithmetic,'' IEEE Std
754-2019.

\bibitem{ahlbach-usatine-pippenger}
C.~Ahlbach, J.~Usatine, and N.~Pippenger, ``Efficient Algorithms for
Zeckendorf Arithmetic.''

\bibitem{dedinechin-flopoco}
F.~de Dinechin and B.~Pasca, ``Designing Custom Arithmetic Data Paths
with FloPoCo,'' \emph{IEEE Design \& Test of Computers}, 2011.

%% PLACEHOLDER: add 4-9 more references as content fills.
%% Drop self-references; drop CLARA / DARPA / IGLA / pellis-vasilev-letter
%% references entirely per anonymisation checklist.

\end{thebibliography}

\end{document}
